\documentclass[11pt]{article}
\usepackage[a4paper,margin=26mm]{geometry}
\usepackage{amsmath,amssymb,amsthm}
\usepackage[T1]{fontenc}
\usepackage{lmodern}
\usepackage{microtype}
\usepackage{hyperref}
\newtheorem{proposition}{Proposition}
\title{An isometric iteration without a weak limit\\\large Counterexample to TLMC 00000008841}
\author{ziangni-sys}
\date{4 October 2026}
\begin{document}
\maketitle
\begin{abstract}
The reflection $T(x)=-x$ on the real Hilbert space is asymptotically
nonexpansive with all coefficients equal to one. It has a fixed point and
preserves a nonempty closed bounded convex interval, yet its Picard orbit
from $1$ does not converge weakly. This disproves the claimed universal
weak convergence of such iterations.
\end{abstract}

\section{Claim and definitions}
The source describes ``Iteration of asymptotically nonexpansive maps'' and
asserts that weak convergence of asymptotically nonexpansive sequences
always holds in Hilbert space. Its further clause concerns strong
convergence and the demiclosedness principle. We disprove the universal
weak-convergence assertion, which suffices to disprove the conjunction.

For a map $T:C\to C$ in a normed space, asymptotic nonexpansiveness means
that there are real numbers $k_n\geq1$, with $k_n\to1$, such that
\[
 \|T^n x-T^n y\|\leq k_n\|x-y\|
 \qquad(x,y\in C,\ n\geq0).
\]
Including $n=0$ does not change the standard definition. The iteration
is the Picard iteration $x_{n+1}=T x_n$, or $x_n=T^n x_0$.
Weak convergence $x_n\rightharpoonup a$ means that
$f(x_n)\to f(a)$ for every continuous linear functional $f$ on the
ambient space. No averaging scheme or asymptotic regularity hypothesis
is stated in the source.

\section{The counterexample}
Let $H=\mathbb R$ with its usual real Hilbert space structure, let
$C=[-1,1]$, and define $T(x)=-x$ on all of $H$. Its restriction maps
$C$ to itself. The interval is nonempty, closed, bounded, and convex;
also $0\in C$ is a fixed point of $T$ and $x_0=1\in C$.

\begin{proposition}
Every iterate of $T$ is an isometry. The orbit $x_n=T^n1$ lies in $C$
and has no weak limit in $H$.
\end{proposition}
\begin{proof}
Negation preserves distance, and a composition of isometries is an
isometry. Induction therefore gives
\[
 |T^n x-T^n y|=|x-y| \qquad(n\geq0).
\]
Consequently $T$ is asymptotically nonexpansive with $k_n=1$ for every
$n$. These same coefficients work for the restriction to $C$.

Since $T(Tx)=x$, induction on $m$ gives the actual orbit identities
\[
 x_{2m}=1,\qquad x_{2m+1}=-1\qquad(m\geq0).
\]
In particular the orbit remains in $C$. If it converged weakly to some
$a\in\mathbb R$, applying the continuous linear identity functional
$f(x)=x$ would give $x_n\to a$ in $\mathbb R$. The even and odd
index maps both tend to infinity, so both subsequences would have limit
$a$. They are constant with respective limits $1$ and $-1$.
Uniqueness of real limits would imply $a=1=-1$, a contradiction.
\end{proof}

This failure occurs despite boundedness, convexity, closedness, and
existence of a fixed point. In particular the demiclosedness clause
cannot restore the already false universal weak-convergence assertion.
The example does not purport to refute results that assume asymptotic
regularity or use an averaged iteration: here
$|x_{n+1}-x_n|=2$ for every $n$.

\section{Lean verification and correspondence}
\raggedright
The accompanying Lean 4.19.0 project uses Mathlib at commit
\texttt{c44e0c8ee63ca166450922a373c7409c5d26b00b}.
In namespace \texttt{AsymptoticIteration8841},
\texttt{reflection} is real negation and \texttt{orbit} is defined by
Lean's actual function iteration, not by a supplied table of values.
\texttt{AsymptoticallyNonexpansive} is the quantified norm inequality
above, expressed with metric distance, and
\texttt{WeaklyConverges} quantifies over all actual continuous real
linear functionals. The use of this definition on $\mathbb R$ is
exactly the standard definition of weak convergence.

\texttt{iterate\_isometry} proves the isometry assertion for every
iterate. \texttt{reflection\_asymptotically\_nonexpansive} supplies
the constant coefficients and their convergence.
\texttt{interval\_geometry} checks the closed, bounded, convex
interval and membership of $0$ and $1$.
\texttt{orbit\_even} and \texttt{orbit\_odd} prove the orbit identities;
the cofinality lemmas and uniqueness of limits then prove
\texttt{orbit\_not\_weakly\_convergent}.
The final theorem \texttt{conjecture\_00000008841} packages the
counterexample, its invariant interval and fixed point, and the
nonexistence of a weak limit.

The proof uses no numerical approximation, unproved axiom, or auxiliary
computation. Its axiom audit reports only the standard logical axioms
\texttt{propext}, \texttt{Classical.choice}, and \texttt{Quot.sound}.
Build with \texttt{lake build}; check the source and axiom audit with
\texttt{lake env lean Main.lean} from the \texttt{lean} directory.
\end{document}
